% A six-cycle: every vertex has a neighbor in each of two coordinates.
\begin{tikzpicture}[
 vertex/.style={draw=BookRule,fill=BookPaper,minimum width=16mm,minimum height=7mm,font=\small},
 first/.style={BookBlue,thick},second/.style={BookAmber,thick,dashed}]
 \node[vertex] (v1) at (0,1.8) {$(a,u)$};
 \node[vertex] (v2) at (1.9,0.9) {$(b,u)$};
 \node[vertex] (v3) at (1.9,-0.9) {$(b,v)$};
 \node[vertex] (v4) at (0,-1.8) {$(c,v)$};
 \node[vertex] (v5) at (-1.9,-0.9) {$(c,w)$};
 \node[vertex] (v6) at (-1.9,0.9) {$(a,w)$};
 \draw[first] (v1)--(v2); \draw[second] (v2)--(v3);
 \draw[first] (v3)--(v4); \draw[second] (v4)--(v5);
 \draw[first] (v5)--(v6); \draw[second] (v6)--(v1);
 \draw[first] (3.2,1.5)--(4,1.5);
 \node[anchor=west,font=\small] at (4.1,1.5) {只改$x_1$的预测};
 \draw[second] (3.2,0.8)--(4,0.8);
 \node[anchor=west,font=\small] at (4.1,0.8) {只改$x_2$的预测};
 \node[anchor=north west,align=left,font=\small,text width=38mm] at (3.2,0.35)
 {每个顶点都连接一条实线和一条虚线，因此DS维为2。\par\smallskip
 找不到两对标签的全部四种组合，因此Natarajan维为1。};
\end{tikzpicture}
